\documentclass[10pt,a4paper]{amsart}
\usepackage[margin=19mm]{geometry}
\usepackage{amsmath,amssymb,mathtools}
\usepackage[scr=rsfs,cal=euler]{mathalfa}
\usepackage{xfrac}
\usepackage[unicode,hidelinks,bookmarks=false]{hyperref}
\newcommand{\ie}{i.e.\ }
\newcommand{\cf}{cf.\ }
\newcommand{\resp}{respectively\ }
\newcommand{\Subsection}{\subsection}
\providecommand{\nobreakdash}{\nobreak\hbox{-}}
\setlength{\emergencystretch}{2em}
\multlinegap0pt
\usepackage[shortlabels]{enumitem}
\usepackage{amssymb}
\usepackage{mathtools}
\usepackage{braket}
\newtheorem{lemma}{Lemma}
\newcommand{\C}{\mathbb{C}}
\newcommand{\R}{\mathbb{R}}
\DeclareMathOperator{\SO}{SO}
\DeclareMathOperator{\spann}{span}
\title{Complete Classification of Directed Quantum Graphs on $M_2$}
\author{Nina Kiefer, Bj\"orn Sch\"afer}
\date{Source: arXiv:2507.15534v2 (CC BY 4.0)}
\begin{document}
\maketitle

\noindent\emph{Source and licence.} Complete Classification of Directed Quantum Graphs on $M_2$. Version arXiv:2507.15534v2 (CC BY 4.0), distributed by arXiv under the Creative Commons Attribution 4.0 International licence, http://creativecommons.org/licenses/by/4.0/, as recorded in the arXiv licence metadata for this exact version and shown on the abstract page https://arxiv.org/abs/2507.15534v2. Full licence notice: \texttt{licenses/arxiv-2507.15534-CC-BY-4.0.txt}.\par\smallskip

\noindent\textit{Editorial scope.} Counted unit: the article's lemma ntg::lemma:2edges\_plane\_classification (source0.tex lines 2951--2984). The article's complete original proof of it is delivered with it (source0.tex lines 2987--3119, one \texttt{proof} environment of 133 lines). No mathematics is written, repaired or re-derived: the statement and the proof are the article's own lines, in the article's own notation, and no retained line is substituted or re-flowed.  Retained uncounted as the source-local context the proof needs: lines 1272--1287; lines 1494--1515.  Nothing was omitted from any retained span, so the package carries no omission record.  No retained line contains a citation, so the wrapper carries no bibliography and no bibliography entry.  No retained line contains a figure environment, an image-inclusion command or drawing code, so no image is delivered and the package declares no asset dependency.  Editorial formatting only, in addition: an AMS \texttt{amsart} preamble that restates the article's own declarations and macros used by the retained lines, namely the environments \texttt{lemma} on separate counters, together with the standard packages those lines need.  The retained environments print in this document as Lemma 1 in order of appearance, whereas the article prints the same items with the numbering of its own full document.  The complete native source of the article as distributed in the arXiv e-print for this exact version is bundled unchanged as \texttt{sources/181947/source0.tex}.
\begin{lemma}
    \label{lin::lemma:existence_of_index_beta_for_lines_in_C3}
        For every non-vanishing vector $v \in \C^3$ there exists some $\lambda \in \C$ and $R \in \SO(3)$ such that one has
        \begin{align*}
            \lambda R v = \begin{pmatrix}
                1 \\ i \beta \\ 0
            \end{pmatrix}
        \end{align*}
        for some $\beta \in [0,1]$. Similarly, for every non-vanishing $v \in \C^2$ there exists some $\lambda \in \C$ and $R \in \SO(2)$ such that
        \begin{align*}
            \lambda R v = \begin{pmatrix}
                1 \\ i \beta
            \end{pmatrix}
        \end{align*}
        holds for some $\beta \in [-1,1]$.
\end{lemma}

\begin{lemma}
    \label{lin::lemma:uniqueness_of_lambda_in_C2}
    For $\beta, \beta^\prime \in [-1,1]$, $\lambda \in \C$ and $R \in \SO(2)$ we have
    \begin{align}
    \label{dir::eq:beta_beta_prime_are_index_of_the_same_line_in_C2}
        \lambda R \begin{pmatrix}
            1 \\ i \beta
        \end{pmatrix}
        = \begin{pmatrix}
            1 \\ i \beta^\prime
        \end{pmatrix}
    \end{align}
    if and only if $\beta = \beta^\prime$ and one of the following is true.
    \begin{enumerate}[label=(\alph*)]
        \item\label{lin::item:beta_beta_prime_are_index_of_the_same_line_in_C2:a} $\beta = \pm 1, \lambda = x + iy \in S^1$ ($x, y \in \R$) and $\displaystyle R = \begin{pmatrix}
            x & -\beta y \\
            \beta y & x \\
        \end{pmatrix}$,
        \item\label{lin::item:beta_beta_prime_are_index_of_the_same_line_in_C2:b} $\beta \in (-1,1), \lambda = \pm 1$ and $\displaystyle R = \lambda \mathrm{I}_2$.
        %\item $\beta \in (-1,1), \lambda = -1$ and $\displaystyle R = -\mathrm{I}_2$.
    \end{enumerate}
\end{lemma}

\begin{lemma}
\label{ntg::lemma:2edges_plane_classification}
    Let $V \subset \C^4$ be a plane with $V \not \perp e_1$. 
    Then there are unique numbers $\beta \in [-1, 1]$ and $\alpha, \gamma, \delta \in \C$ such that 
    \begin{align*}
        \left[\begin{array}{c|c|c}
            1& 0 & 0 \\
            \hline
            0 & R & 0 \\
            \hline 
            0 & 0 & 1
            \end{array}\right]
        V = 
            \spann \left\{ \begin{pmatrix}
                0 \\ 1 \\ i \beta \\ \delta
            \end{pmatrix}, \begin{pmatrix}
                1 \\ i \alpha \beta - \gamma \overline{\delta} \\ \alpha \\ \gamma
            \end{pmatrix} \right\}
            =\vcentcolon V^{(q,2)}_{\alpha,\beta,\gamma,\delta}
    \end{align*}
    holds for some $R \in \SO(2)$ and
    \begin{align}
        \label{8::eq:J^(q,2)}\tag{$\ast$}
        \begin{aligned}
            \beta = \pm 1, \delta \neq 0 &\implies \delta \in (0, \infty), \\
            \beta = \pm 1, \delta = 0 &\implies \alpha \in [0, \infty), \\
            \beta \in (-1,1), \delta \neq 0 &\implies \mathrm{arg}(\delta) \in [0, \pi), \\
            \beta \in (-1,1), \delta = 0 &\implies \mathrm{arg}(\alpha) \in [0, \pi).
        \end{aligned}
    \end{align}
    Denote the index set 
    $J^{(q,2)}\vcentcolon=
        \{(\alpha,\beta,\gamma,\delta)\vcentcolon\text{\eqref{8::eq:J^(q,2)} is fulfilled}\}$.
\end{lemma}

\begin{proof}
    Up to scalar multiplication there is a unique non-vanishing vector $v \in V$ of the form
    \begin{align*}
        v = \begin{pmatrix}
            0 \\ v_1 \\ v_2 \\ v_3
        \end{pmatrix} \in V.
    \end{align*}
    Indeed, if there were two linearly independent vectors with vanishing first component, they would span the whole space $V$ -- in contradiction to $V \not \perp e_1$. Using Lemma \ref{lin::lemma:existence_of_index_beta_for_lines_in_C3} one finds a unique $\beta \in [-1,1]$ such that
    \begin{align*}
        \begin{pmatrix}
            v_1 \\ v_2
        \end{pmatrix}
        = 
        \lambda R \begin{pmatrix}
                1 \\ i \beta
        \end{pmatrix}
    \end{align*}
    holds for some $\lambda \in \C$ and $R \in \SO(2)$. Consequently, we have with $\delta := \lambda v_3$
    \begin{align*}
        \begin{pmatrix}
                0 \\ 1 \\ i \beta \\ \delta
            \end{pmatrix}
        = 
        \lambda \left[\begin{array}{c|c|c}
            1& 0 & 0 \\
            \hline
            0 & R & 0 \\
            \hline 
            0 & 0 & 1
            \end{array}\right]
            v.
    \end{align*}
    Let $w \in V$ be another vector that is orthogonal to $v$ and such that $V = \spann \{ v, w\}$. Then $w$ has non-vanishing first component since $v$ is not orthogonal to $e_1$ and without loss of generality we may assume that its first component is $w_1 = 1$. Using $v \perp w$ one readily checks that there are numbers $\alpha, \gamma \in \C$ such that
    \begin{align*}
        \begin{pmatrix}
                1 \\ i \alpha \beta - \gamma \overline{\delta} \\ \alpha \\ \gamma
            \end{pmatrix}
        =
        \left[\begin{array}{c|c|c}
            1& 0 & 0 \\
            \hline
            0 & R & 0 \\
            \hline 
            0 & 0 & 1
            \end{array}\right]
            w.
    \end{align*}
    It remains to investigate uniqueness of the numbers $\alpha, \beta, \gamma, \delta \in \C$. 
    For that, assume
    \begin{align*}
        \spann \left\{ \begin{pmatrix}
            0 \\ 1 \\ i \beta \\ \delta
        \end{pmatrix}, \begin{pmatrix}
            1 \\ i \alpha \beta - \gamma \overline{\delta} \\ \alpha \\ \gamma
        \end{pmatrix} \right\}
        = 
        \left[\begin{array}{c|c|c}
            1& 0 & 0 \\
            \hline
            0 & R & 0 \\
            \hline 
            0 & 0 & 1
            \end{array}\right]
            \spann \left\{ \begin{pmatrix}
                0 \\ 1 \\ i \beta^\prime \\ \delta^\prime
            \end{pmatrix}, \begin{pmatrix}
                1 \\ i \alpha^\prime \beta^\prime - \gamma^\prime \overline{\delta^\prime} \\ \alpha^\prime \\ \gamma^\prime
            \end{pmatrix} \right\}
    \end{align*}
    for some $R \in \SO(2)$ and $\beta^\prime \in [-1,1]$, $\alpha^\prime, \gamma^\prime, \delta^\prime \in \C$. Looking at the first components of the spanning vectors one readily checks
    \begin{align}
        \begin{pmatrix}
            0 \\ 1 \\ i \beta \\ \delta
        \end{pmatrix}
        &=
        \lambda' \left[\begin{array}{c|c|c}
            1& 0 & 0 \\
            \hline
            0 & R & 0 \\
            \hline 
            0 & 0 & 1
        \end{array}\right]
        \begin{pmatrix}
            0 \\ 1 \\ i \beta^\prime \\ \delta^\prime
        \end{pmatrix} 
    \end{align}
    for some $\lambda' \in \C$ and
    \begin{align}
    \label{ntg::eq:plane_lemma_case3}
        \begin{pmatrix}
            1 \\ i \alpha \beta - \gamma \overline{\delta} \\ \alpha \\ \gamma
        \end{pmatrix}
        &= 
        \left[\begin{array}{c|c|c}
            1& 0 & 0 \\
            \hline
            0 & R & 0 \\
            \hline 
            0 & 0 & 1
        \end{array}\right]
        \begin{pmatrix}
            1 \\ i \alpha^\prime \beta^\prime - \gamma^\prime \overline{\delta^\prime} \\ \alpha^\prime \\ \gamma^\prime
        \end{pmatrix}.
    \end{align}
    The first equation entails $\delta = \lambda' \delta^\prime$ as well as $\beta = \beta^\prime$ (due to Lemma \ref{lin::lemma:uniqueness_of_lambda_in_C2}), while the second entails $\gamma = \gamma^\prime$. 
    Thus, the numbers $\beta, \gamma$ are unique. 
    By Lemma \ref{lin::lemma:uniqueness_of_lambda_in_C2} the two equations hold if and only if -- additional to $\beta = \beta^\prime$, $\gamma = \gamma^\prime$ and $\delta = \lambda'\delta^\prime$ -- one of the following is true.
    \begin{enumerate}[label=(\alph*)]
        \item $\beta =\pm 1$ and $\displaystyle R = \begin{pmatrix}
            x & -\beta y \\
            \beta y & x
        \end{pmatrix}$
        for $x+ iy = \lambda' \in S^1$. By choosing $\lambda'$ appropriately we can obtain $\delta \in [0, \infty)$. If we had $\delta, \delta^\prime \in (0, \infty)$, then one checks immediately $\lambda' = 1$ and thus, $R=\mathrm{I}_2$ and $\alpha = \alpha^\prime$. Thus, in this case, it  follows that the numbers $\alpha, \beta, \gamma, \delta$ are unique if one asks $\delta \in [0, \infty)$. If, on the other hand, $\delta = 0$, then (\ref{ntg::eq:plane_lemma_case3}) is equivalent to
        \begin{align*}
            \alpha \begin{pmatrix}
                i \beta \\ 1
            \end{pmatrix}
            = \alpha^\prime 
            \begin{pmatrix}
                x & -\beta y \\
                \beta y & x
            \end{pmatrix}
            \begin{pmatrix}
                i \beta^\prime \\ 1
            \end{pmatrix}
        \end{align*}
        A simple calculation shows that this is true if and only if $\alpha = \lambda' \alpha^\prime$. 
        %or $\alpha = \overline{\lambda'} \alpha^\prime$, respectively, depending on whether $\beta = 1$ or $\beta = -1$. 
        Thus, in that case one can choose $\alpha \in [0, \infty)$, and the numbers $\alpha, \beta, \gamma, \delta$ are unique if one asks $\alpha \in [0, \infty)$.
        \item $\beta \in (-1,1)$, $\lambda' = \pm 1$ and $R =\lambda^\prime \mathrm{I}_2$. In this case, the two equations entail $\alpha = \pm\alpha^\prime$ and $\delta = \pm\delta^\prime$.
    \end{enumerate}
    The latter two cases show that the pair $(\alpha, \delta)$ is unique up to a sign if $\beta \in (-1,1)$. Altogether, the statement follows.
\end{proof}

\end{document}
